\documentclass[twoside]{article}
\usepackage[utf8]{inputenc}

% The package supports other formats as well, but this template
% is primarily checked and updated for examinations at TUM BGDM
%\usepackage[uebung, answers]{tumbgdm}

\usepackage[uebung]{tumbgdm}
\usepackage{tabularx}
\usepackage{amsmath}

\Professor{Prof. Dr. Martin Werner}
\Universitaet{Technical University of Munich}
\Institut{School of Engineering and Design}
\Lehrstuhl{Professorship Big Geospatial Data Management}
\Semester[WiTe 2024/25]{Winter Term 2024/25}
\Vorlesung{Principles of Programming}
\InstitutsWebsite{https://www.bgd.lrg.tum.de/}
\VorlesungsWebsite{https://www.moodle.tum.de/course/view.php?id=103704}

\nummerUebungsblatt{10}
\nummerErsteFrage{1}

\LearningOutcome{ Multilateration}

\begin{document}
\maketitle

\vspace{10pt}

"Pseudo-range multilateration, often simply multilateration when in context, is a technique for determining the position of an unknown point, such as a vehicle, based on measurement of the times of arrival of energy waves traveling between the unknown point and multiple stations at known locations." [\href{https://en.wikipedia.org/wiki/Pseudo-range\_multilateration}{Wikipedia}]

\begin{task}{Multilateration}{}{}

Write a program to perform multilateration.
You should first review the lecture slides and familiarize yourself with the algorithm.
\vspace{10pt}

\noindent Following is a summary of the algorithm:
\begin{enumerate}
    \item Set up a non-linear system of equations from measurements.
    \item Create a linearized version locally by choosing an initial value and using Taylor expansion at this point.
    \item Solve this equation and obtain an update vector.
    \item Update this and go back to b. with this new place information
\end{enumerate}
\vspace{5pt}
Thus, implement the following steps:
\begin{enumerate}
    \item Calculate distances:
$\alpha_{i}=\sqrt{(x_{i}-\Tilde{x})^2+(y_{i}-\Tilde{y})^2}$ for $i = 1, ..., 4$.
    \item Build matrix:
$A_{i}=(-\frac{(x_{i}-\Tilde{x})}{\alpha_{i}}, -\frac{(y_{i}-\Tilde{y})}{\alpha_{i}})$.
    \item Create: $b_{i}=d_{i}-\alpha_{i}$.
    \item Solve: $A\begin{pmatrix}\hat{x}\\\hat{y}\end{pmatrix}=b$.
    \item Compute: $\begin{pmatrix}\Tilde{x}_{2}\\\Tilde{y}_{2}\end{pmatrix}= \begin{pmatrix}\Tilde{x}_{1}\\\Tilde{y}_{1}\end{pmatrix} + \begin{pmatrix}\hat{x}\\\hat{y}\end{pmatrix}$.
\end{enumerate}

\clearpage
You can use the following data from the lecture:\\
\\
\begin{tabularx}{0.8\textwidth} { 
  | >{\raggedright\arraybackslash}X 
  | >{\raggedright\arraybackslash}X 
  | >{\raggedright\arraybackslash}X | }
 \hline
 \textbf{Component} & \textbf{Coordinate} & \textbf{Measurement} ($d_{j}$) \\
 \hline
 Base station 1  & $(x_{1}, y_{1}) = (0, 0)$  & $2.92$ \\
 \hline
 Base station 2  & $(x_{2}, y_{2}) = (10, 0)$  & $8.14$ \\
 \hline
 Base station 3  & $(x_{3}, y_{3}) = (15, 10)$  & $15.46$ \\
 \hline
 Base station 4  & $(x_{4}, y_{4}) = (0, 12)$  & $9.89$ \\
 \hline
 Mobile Device  & $(x_{e}, y_{e}) = (2, 2)$  & $-$ \\
 \hline
 Initial Location Estimate  & $(\Tilde{x}, \Tilde{y}) = (20, 20)$  & $-$ \\
\hline
\end{tabularx}

\end{task}


\end{document}